\subsection{Jets of sections of a vector bundle}

12.6.1. Let E be a vector bundle of class \(C^{r}\) with base X, and let \(\pi\) be its projection. Let \(c_{0} = (\mathrm{U}, \varphi, \mathrm{F}_{0})\) be a vector chart of E and \(c_{1} = (\mathrm{U}, \psi, \mathrm{F}_{1})\) a chart of the manifold X, with the same domain U. These charts define a bijection \(\theta\) of \(\mathrm{P}^{k}(\pi)|\mathrm{U}\) onto \(\mathrm{J}^{k}(\psi(\mathrm{U}), \mathrm{F}_{0}) = \psi(\mathrm{U}) \times \mathrm{F}_{0} \times \mathrm{Q}^{k}(\mathrm{F}_{1}, \mathrm{F}_{0})\), cf. 12.2.2 and 12.3.1, whence a vector chart

\[
d = (\mathrm{U}, \theta, \mathrm{G}), \quad \text{with}\quad \mathrm{G} = \mathrm{F}_{0}\times\mathrm{Q}^{k}(\mathrm{F}_{1},\mathrm{F}_{0}) = \prod_{m=0}^{k}\mathrm{P}_{m}(\mathrm{F}_{1};\mathrm{F}_{0})
\]

of \(\mathbf{P}^{k}(\pi)\). The charts thus obtained form a \(\mathbf{C}^{r - k}\)-atlas, which endows \(\mathbf{P}^{k}(\pi)\) with a vector-bundle structure of class \(\mathbf{C}^{r - k}\), with base X. This vector bundle is denoted \(\mathbf{P}^{k}(\mathbf{E})\); the underlying manifold structure is the one defined in 12.5.1.

Let U be an open subset of X, and let \(f \in \mathcal{S}_{\mathrm{E}}^{r}(\mathrm{U})\) be a section of class \(\mathbf{C}^{r}\) of E over U. The map \(j^{k}(f): x \mapsto j_{x}^{k}(f)\) is a section of class \(\mathbf{C}^{r - k}\) of \(\mathrm{P}^{k}(\mathrm{E})\) over U. The map \(j^{k}: \mathcal{S}_{\mathrm{E}}^{r}(\mathrm{U}) \to \mathcal{S}_{\mathrm{P}^{k}(\mathrm{E})}^{r - k}(\mathrm{U})\) is K-linear.

12.6.2. If F is a Banach space, the identification (12.5.2) of \(\mathrm{P}^{k}(\mathrm{F}_{\mathbf{x}})\) with \(\mathrm{J}^{k}(\mathrm{X},\mathrm{F})\) endows the latter manifold with a vector-bundle structure of class \(C^{r-k}\), with base X. When F = K, one writes \(\mathrm{P}^{k}(\mathrm{X})\) instead of \(\mathrm{P}^{k}(\mathrm{K}_{\mathbf{x}})\).

Example. --- Take \(X = K^{n}\), and denote by \(u_{1}, \ldots, u_{n}\) the coordinate functions on \(K^{n}\); then the fiber \(\mathrm{P}_{0}^{k}(\mathrm{X})\) of \(\mathrm{P}^{k}(\mathrm{X})\) at 0 has as basis the family of jets of order k

of the monomials \(u_{1}^{m_1} \ldots u_{n}^{m_n}\), with \(m_i \geqslant 0\), \(\sum_{i=0}^{n} m_i \leqslant k\).

12.6.3. Let \(d\) be an integer \(\geqslant 0\) , \(\mathrm{E}_{1},\ldots,\mathrm{E}_{d}\) , \(\mathbf{F}\) vector bundles of class \(\mathbf{C}^{r}\) with base \(\mathbf{X}\) , and let \(u\colon\mathrm{E}_{1}\times_{\mathbf{x}}\cdots\times_{\mathbf{x}}\mathrm{E}_{d}\to\mathbf{F}\) be a multilinear morphism (7.3.1) of class \(\mathbf{C}^{r}\) . Then there exists a unique multilinear morphism

\[
\mathrm{P} ^ {k} (u) \colon \mathrm{P} ^ {k} (\mathrm{E} _ {1}) \times_ {\mathrm{x}} \dots \times_ {\mathrm{x}} \mathrm{P} ^ {k} (\mathrm{E} _ {d}) \rightarrow \mathrm{P} ^ {k} (\mathrm{F})
\]

such that

\[
\mathrm{P} ^ {k} (u) \left(j ^ {k} \left(s _ {1}\right), \dots , j ^ {k} \left(s _ {d}\right)\right) = j ^ {k} \left(u \left(s _ {1}, \dots , s _ {d}\right)\right)
\]

for every open subset U of X and every sequence of sections \(s_i \in \mathcal{S}_{\mathbf{E}_i}^r(\mathbf{U})\), \(1 \leqslant i \leqslant d\).

If A is a bundle of algebras (7.3.2) with base X, the morphism from \(\mathbf{P}^{k}(\mathbf{A}) \times_{\mathbf{x}} \mathbf{P}^{k}(\mathbf{A})\) to \(\mathbf{P}^{k}(\mathbf{A})\) deduced from the multiplication \(A \times_{x} A \to A\) makes \(\mathbf{P}^{k}(\mathbf{A})\) into a bundle of algebras; if A is a bundle of associative algebras, \(\mathbf{P}^{k}(\mathbf{A})\) is a bundle of associative algebras; if moreover M is a bundle of A-modules (7.3.3), \(\mathbf{P}^{k}(\mathbf{M})\) is a bundle of \(\mathbf{P}^{k}(\mathbf{A})\)-modules. In particular, \(\mathbf{P}^{k}(\mathbf{X})\) is a bundle of associative, commutative, and unital algebras; if E is a vector bundle, \(\mathbf{P}^{k}(\mathbf{E})\) is a bundle of \(\mathbf{P}^{k}(\mathbf{X})\)-modules. If \(E \to F\) is a morphism of vector bundles, then \(\mathbf{P}^{k}(u): \mathbf{P}^{k}(\mathbf{E}) \to \mathbf{P}^{k}(\mathbf{F})\) is a \(\mathbf{P}^{k}(\mathbf{X})\)-homomorphism.

12.6.4. If \(E \xrightarrow{u} F \xrightarrow{v} G\) is a locally split exact sequence of vector bundles with base X, then so is the sequence

\[
\mathbf {P} ^ {k} (\mathrm{E}) \rightarrow \mathbf {P} ^ {k} (\mathrm{F}) \rightarrow \mathbf {P} ^ {k} (\mathrm{G})
\]

where the homomorphisms considered are \(\mathrm{P}^{k}(u)\) and \(\mathrm{P}^{k}(v)\).

12.6.5. Let \(k'\) and \(k''\) be two positive integers with sum \(k\), and let E be a vector bundle of class \(C^{r}\) and base X. The morphism

\[
\alpha \colon \mathrm{J} ^ {k} (\mathrm{X}, \mathrm{E}) \rightarrow \mathrm{J} ^ {k ^ {\prime \prime}} (\mathrm{X}, \mathrm{J} ^ {k ^ {\prime}} (\mathrm{X}, \mathrm{E})) \quad (\text { cf. } 1 2. 3. 8)
\]

induces a morphism of vector bundles

\[
\beta \colon \mathbf {P} ^ {k} (\mathrm{E}) \rightarrow \mathbf {P} ^ {k ^ {\prime \prime}} (\mathbf {P} ^ {k ^ {\prime}} (\mathrm{E}))
\]

which is of class \(C^{r-k}\). If U is an open subset of X and \(f\in\mathcal{S}_{\mathrm{E}}^{r}(\mathrm{U})\), one has

\[
\beta (j ^ {k} (f)) = j ^ {k ^ {\prime \prime}} (j ^ {k ^ {\prime}} (f)).
\]

When K is of characteristic zero, \(\beta\) is an isomorphism of \(\mathbf{P}^{k}(\mathbf{E})\) onto a vector subbundle of \(\mathbf{P}^{k^{\prime\prime}}(\mathbf{P}^{k^{\prime}}(\mathbf{E}))\).

12.6.6. Let \(k'\) be an integer such that \(0 \leqslant k' \leqslant k\), and let E be a vector bundle of class \(C^r\) and base X. The map

\[
r ^ {k, k ^ {\prime}}: \mathrm{P} ^ {k} (\mathrm{E}) \rightarrow \mathrm{P} ^ {k ^ {\prime}} (\mathrm{E})
\]

is a locally split surjective morphism of class \(C^{r-k}\). Its kernel \(\mathbf{N}^{k,k'}(\mathbf{E})\) consists of the jets of sections of E having contact of order \(\geqslant k'\) with the zero section.

12.6.7 ("The vector functor \(P_{m}\)"). Using the notation of 7.6, 7.7, and 7.8, set \(I_{+} = \{0\}\), \(I_{-} = \{1\}\), and let m be an integer \(\geqslant 0\). If \(\mathcal{V} = (\mathrm{V}_{0}, \mathrm{V}_{1})\) is a pair of Banach spaces, denote by \(\tau_{m}(\mathcal{V})\) the Banach space \(\mathrm{P}_{m}(\mathrm{V}_{1}; \mathrm{V}_{0})\) of continuous homogeneous polynomials of degree m on \(V_{1}\) with values in \(V_{0}\) (A.2). Likewise, if \(f = (f_{0}, f_{1})\), where \(f_{0}: V_{0} \to V_{0}'\) and \(f_{1}: V_{1}' \to V_{1}\) are morphisms of Banach spaces, denote by \(\tau_{m}(f)\) the morphism \(p \mapsto f_{0} \circ p \circ f_{1}\) from \(\mathrm{P}_{m}(\mathrm{V}_{1}; \mathrm{V}_{0})\) to \(\mathrm{P}_{m}(\mathrm{V}_{1}', \mathrm{V}_{0}')\). One thus obtains a vector functor \(\tau_{m}\) of class \(\mathbf{C}^{\omega}\). If \(\mathbf{E}_{0}\) and \(\mathbf{E}_{1}\) are two vector bundles with base \(\mathbf{X}\), one denotes by \(\mathbf{P}_{m}(\mathbf{E}_{1};\mathbf{E}_{0})\) the vector bundle deduced from \((\mathbf{E}_{0},\mathbf{E}_{1})\) by \(\tau_{m}\) (7.6.2).

12.6.8. Let us resume the notations and hypotheses of 12.6.1. There exists one and only one morphism of vector bundles

\[
\iota\colon\mathrm{P}_{k}(\mathrm{T}(\mathrm{X});\mathrm{E})\rightarrow\mathrm{P}^{k}(\mathrm{E})\quad(\text{cf. 12.6.6}),
\]

such that, whatever the vector chart \(c_{0} = (\mathrm{U}, \varphi, \mathrm{F}_{0})\) of E and the chart \(c_{1} = (\mathrm{U}, \psi, \mathrm{F}_{1})\) of X, the following diagram is commutative:

\[
\begin{array}{c c c} \mathrm{P} _ {k} (\mathrm{T} (\mathrm{X}), \mathrm{E}) | \mathrm{U} & \xrightarrow {\iota | \mathrm{U}} & \mathrm{P} ^ {k} (\mathrm{E}) | \mathrm{U} \\ \Biggl \downarrow_ {\eta} & & \Biggl \downarrow_ {\theta} \\ \mathrm{U} \times \mathrm{P} _ {k} (\mathrm{F} _ {1}; \mathrm{F} _ {0}) & \xrightarrow {i} & \mathrm{U} \times \prod_ {m = 0} ^ {k} \mathrm{P} _ {m} (\mathrm{F} _ {1}; \mathrm{F} _ {0}) \end{array}
\]

where: 1) \(\iota|_{\mathbf{U}}\) is the restriction of \(\iota\) to \(\mathrm{P}_{k}(\mathrm{T}(\mathbf{X}),\mathbf{E})|_{\mathbf{U}}\);

\begin{enumerate}
\def\labelenumi{\arabic{enumi})}
\setcounter{enumi}{1}
\item
  \(\theta\) is the bijection defined in 12.6.1;
\item
  i is the map \((u, p) \mapsto (u, 0, \ldots, 0, p)\);
\item
  \(\eta\) is deduced, by means of the vector functor \(\tau_{k}\) of the vector chart \(c_{1}^{\prime}\) of T(X) (cf. 8.1.1) and of the vector chart \(c_{0}\) of E.
\end{enumerate}

The morphism \(\iota\) is an isomorphism of \(\mathrm{P}_{k}(\mathrm{T}(\mathrm{X});\mathrm{E})\) onto the vector subbundle \(\mathbf{N}^{k,k - 1}(\mathbf{E})\) of \(\mathrm{P}^k (\mathbf{E})\) (cf. 12.6.6); the sequence

\[
0 \rightarrow \mathrm{P}_{k}(\mathrm{T}(\mathrm{X});\mathrm{E}) \xrightarrow {\iota} \mathrm{P}^{k}(\mathrm{E}) \xrightarrow {r^{k,k-1}} \mathrm{P}^{k-1}(\mathrm{E}) \rightarrow 0
\]

is a locally split exact sequence of vector bundles.

More generally, let us set \(\mathrm{N}_{0}=\mathrm{P}^{k}(\mathrm{E})\), \(\mathrm{N}_{m}=\mathrm{N}^{k,m-1}(\mathrm{E})\) for \(m\geqslant1\), so that the \(N_{m}\) form a decreasing sequence of vector subbundles of \(\mathrm{P}^{k}(\mathrm{E})\):

\[
\mathrm{P} ^ {k} (\mathrm{E}) = \mathrm{N} _ {0} \supset \mathrm{N} _ {1} \supset \dots \supset \mathrm{N} _ {k} \supset \mathrm{N} _ {k + 1} = 0.
\]

For \(0 \leqslant m \leqslant k\), the projection \(r^{k,m}: \mathbf{P}^{k}(\mathbf{E}) \to \mathbf{P}^{m}(\mathbf{E})\) defines an isomorphism of \(\mathbf{N}_{m}/\mathbf{N}_{m+1}\) onto \(\mathbf{N}^{m,m-1}(\mathbf{E})\); in view of what precedes, one thus obtains isomorphisms

\[
\iota_ {m} \colon \mathrm{N} _ {m} / \mathrm{N} _ {m + 1} \rightarrow \mathrm{P} _ {m} (\mathrm{T} (\mathrm{X}); \mathrm{E}) \quad (0 \leqslant m \leqslant k).
\]

In particular, one has \(N_{0}/N_{1} \simeq E\) and \(N_{1}/N_{2} \simeq \mathcal{L}(T(X); E)\) . For k = 1, this gives an exact sequence:

\[
0 \rightarrow \mathscr {L} (\mathrm{T} (X); E) \rightarrow \mathrm{P} ^ {1} (E) \rightarrow E \rightarrow 0.
\]

